\documentclass[11pt]{article}
\usepackage[margin=1in]{geometry}
\usepackage[T1]{fontenc}
\usepackage{amsmath,amssymb,amsthm,booktabs}
\usepackage[colorlinks=true,linkcolor=blue,urlcolor=blue]{hyperref}
\newtheorem{theorem}{Theorem}[section]
\newtheorem{proposition}[theorem]{Proposition}
\newtheorem{lemma}[theorem]{Lemma}
\newtheorem{corollary}[theorem]{Corollary}
\theoremstyle{definition}
\newtheorem{definition}[theorem]{Definition}
\newtheorem{example}[theorem]{Example}
\theoremstyle{remark}
\newtheorem{remark}[theorem]{Remark}
\newcommand{\Pset}{\mathcal P}
\newcommand{\leg}[2]{\left(\frac{#1}{#2}\right)}
\newcommand{\ind}{\mathbf 1}
\title{Local sieves, correlations, and witness counts\\for shifted even squares}
\author{Mohammad Aldebbeh}
\date{30 September 2026}
\begin{document}
\maketitle
\begin{abstract}
For the shifted even square $G_a(m)=4m^2+a$, we count residue classes
avoiding a finite set of prime divisors or prime-square divisors.
Quadratic residues and the Chinese remainder theorem give exact
densities, including joint densities for several shifts at the same
center. We classify the roots modulo $p^2$ by an elementary lifting
argument and obtain explicit local counts for the opposite-sign pair
$4m^2\pm q$, with $q$ an odd prime. This pair has no local obstruction
for either finite sieve. Exact computations implement the formulas;
a separate factorization experiment compares prime witnesses with
rough squarefree almost-prime witnesses. Finite moment inequalities
relate the number of witnesses to the number of occupied centers.
The results concern finite local sieves and do not establish global
squarefreeness or simultaneous prime values.
\end{abstract}
\tableofcontents

\section{Problem and notation}
For an even positive integer $n=2m$ and an odd prime $q$, consider
\[
 n^2+q=4m^2+q,\qquad n^2-q=4m^2-q.
\]
Before asking whether these values are prime, one can ask which centers
are ruled out by a particular divisor. That question is periodic and
can be answered exactly. The same idea works for prime squares and
provides a finite local test towards squarefreeness.

\begin{definition}
For an integer $a$, put $G_a(m)=4m^2+a$. For $d\ge1$, define
\[
 R_a(d)=\{r\in\{0,\ldots,d-1\}:d\mid G_a(r)\},
 \qquad \rho_a(d)=|R_a(d)|.
\]
Let $\Pset$ be a finite set of distinct primes, and write
\[
 P=\prod_{p\in\Pset}p,\qquad M=\prod_{p\in\Pset}p^2=P^2.
\]
For a finite set of integer offsets $\mathcal A$, a prime-sieve survivor
is a center for which $p\nmid G_a(m)$ for all $a\in\mathcal A$ and
$p\in\Pset$. A prime-square-sieve survivor instead requires
$p^2\nmid G_a(m)$ for every such $a,p$.
\end{definition}

All products in this note are finite. An empty product is $1$;
an empty offset family imposes no restrictions. Prime-sieve survival
has period $P$, and prime-square-sieve survival has period $M$.
These are valid periods, not necessarily the smallest ones.
The local statements allow negative centers and arbitrary integer
offsets, including zero.

There are two counting questions in the project. Fixing offsets and
varying $m$ gives a local sieve count. Fixing $n$ and searching over
prime offsets $q$ gives a witness count. The experiments measure both,
but the two counts have different quantifiers. In particular, saying
that both signs have a witness allows different offsets on the two
sides; asking for the same $q$ to work twice is stronger.

\begin{definition}
For $v>1$, write $v=\prod_p p^{e_p}$ and $\Omega(v)=\sum_p e_p$.
A prime witness is a prime value $v=n^2\pm q$. A rough squarefree
$P_2$ witness is a value $v>1$ with every $e_p\le1$,
$\Omega(v)\le2$, and least prime divisor $P^-(v)>n^{1/4}$.
\end{definition}

Here $P_2$ means at most two prime factors, counted with multiplicity.
Primes can qualify, whereas prime squares cannot. The program tests
roughness by $(P^-(v))^4>n$, with exact integers. Values at most $1$
are in neither witness class. These are conditions on the fully
factored value, stronger than passing a fixed local sieve.

For an odd prime $p$, the Legendre symbol $\leg{t}{p}$ is $0$ if
$p\mid t$, $1$ if $t$ is a nonzero square modulo $p$, and $-1$
otherwise. This definition is enough for the root counts below.

\section{Local divisibility modulo primes}
\begin{proposition}[Prime root count]\label{prop:roots}
For an integer $a$ and an odd prime $p$,
\[
 p\mid G_a(m)\iff (2m)^2\equiv-a\pmod p,
 \qquad \rho_a(p)=1+\leg{-a}{p}.
\]
At $2$, the count is $0$ for odd $a$ and $2$ for even $a$.
\end{proposition}
\begin{proof}
Multiplication by $2$ is a bijection modulo an odd prime.
The equation $x^2=0$ has just the root $0$ in the field $\mathbb F_p$.
A nonzero square has two distinct roots $x,-x$: any other root $y$
would satisfy $(y-x)(y+x)=0$, so $y=x$ or $y=-x$.
A nonsquare has no root by definition. Finally,
$4m^2+a\equiv a\pmod2$ for both classes of $m$.
\end{proof}

The usual description ``zero or two roots'' therefore requires
$p\nmid a$. For $a=\pm q$ and $p=q$, the single bad class is
$m\equiv0\pmod q$. At $2$, neither odd offset removes a center.

\begin{proposition}[Squarefree moduli]\label{prop:crt}
For squarefree $d\ge1$,
\[
 \rho_a(d)=\prod_{p\mid d}\rho_a(p).
\]
If $a$ is odd, this is at most $2^{\omega(d)}$, where $\omega(d)$
counts the distinct prime factors. At $d=1$, the count is $1$.
\end{proposition}
\begin{proof}
The Chinese remainder theorem (CRT) bijects classes modulo $d$ with
tuples of classes modulo its prime factors. Under this bijection,
the roots form the Cartesian product of the local root sets.
Taking cardinalities proves the formula and the preceding proposition
gives the bound. Modulo $1$, the sole canonical class is $0$.
\end{proof}

CRT combines coprime moduli. It does not treat repeated prime factors
as independent, or different offsets at one prime as independent.

\section{Finite prime sieves and interval counts}
For integer endpoints $L\le U$, put $H=U-L$ and
\[
 A_a(d;L,U)=\#\{m:L<m\le U,\ d\mid G_a(m)\}.
\]
\begin{lemma}[Count in one residue class]\label{lem:interval}
For $d\ge1$ and $0\le r<d$, the number of integers in $(L,U]$
congruent to $r$ modulo $d$ is
\[
 \left\lfloor\frac{U-r}{d}\right\rfloor-
 \left\lfloor\frac{L-r}{d}\right\rfloor.
\]
It differs from $H/d$ by at most $1$. Hence
\[
 \left|A_a(d;L,U)-H\frac{\rho_a(d)}d\right|\le\rho_a(d).
\]
\end{lemma}
\begin{proof}
Write an integer in the class as $r+kd$ and count $k$ between the
two endpoints. Each floor differs from its argument by a number in
$(-1,0]$, giving the error bound. Sum over the roots in $R_a(d)$.
\end{proof}

Write $S_a(L,U;P)$ for the prime-sieve survivor count of one shift,
and put $\delta_a(P)=\prod_{p\mid P}(1-\rho_a(p)/p)$.
\begin{theorem}[Finite inclusion-exclusion]\label{thm:sieve}
For $P$ a product of distinct primes,
\[
 S_a(L,U;P)=\sum_{d\mid P}\mu(d)A_a(d;L,U),
 \qquad \mu(d)=(-1)^{\omega(d)}.
\]
Moreover,
\[
 \left|S_a(L,U;P)-H\delta_a(P)\right|
 \le\prod_{p\mid P}(1+\rho_a(p)).
\]
For odd $a$ and $k$ odd sieving primes, the bound is at most $3^k$.
\end{theorem}
\begin{proof}
For each $m$, expand the finite product
\[
 \ind_{\gcd(G_a(m),P)=1}
 =\prod_{p\mid P}(1-\ind_{p\mid G_a(m)})
 =\sum_{d\mid P}\mu(d)\ind_{d\mid G_a(m)}.
\]
Summing gives the identity. Lemma~\ref{lem:interval} bounds the total
error by $\sum_{d\mid P}\rho_a(d)$. The CRT factorization gives
\[
 \sum_{d\mid P}\rho_a(d)=\prod_{p\mid P}(1+\rho_a(p)),
 \qquad
 \sum_{d\mid P}\frac{\mu(d)\rho_a(d)}d
 =\prod_{p\mid P}\left(1-\frac{\rho_a(p)}p\right).
\]
For odd $a$, the factor at $2$ is $1$, and each other error factor
is at most $3$.
\end{proof}

Adding primes excludes more divisors, but this elementary error bound
grows. It can exceed the interval length. For short windows, the exact
residue-class count will be more useful than replacing it by a density.

\section{Several offsets at the same center}
\begin{lemma}[Overlap of bad classes]\label{lem:overlap}
For integers $a,b$ and any positive modulus $d$, the sets $R_a(d)$
and $R_b(d)$ are identical if $a\equiv b\pmod d$, and disjoint otherwise.
\end{lemma}
\begin{proof}
Congruent offsets define the same congruence. A common root would
give $d\mid G_a(m)-G_b(m)=a-b$.
\end{proof}

\begin{proposition}[Local covariance]\label{prop:covariance}
Choose $m$ uniformly modulo a prime $p$ and let $I_a=\ind_{p\mid G_a(m)}$.
Then
\[
 \operatorname{Cov}(I_a,I_b)=
 \begin{cases}
 \displaystyle\frac{\rho_a(p)}p\left(1-\frac{\rho_a(p)}p\right),
    &a\equiv b\pmod p,\\[6pt]
 \displaystyle-\frac{\rho_a(p)\rho_b(p)}{p^2},&a\not\equiv b\pmod p.
 \end{cases}
\]
\end{proposition}
\begin{proof}
The expectations of $I_a,I_b$ are $\rho_a(p)/p,\rho_b(p)/p$.
By Lemma~\ref{lem:overlap}, the expectation of $I_aI_b$ is
$\rho_a(p)/p$ for congruent offsets and zero otherwise.
Subtract the product of the expectations; in the first case the
two root counts agree.
\end{proof}

For $a=5,b=11,p=3$, both root sets are $\{1,2\}$ and the covariance
is $2/9$. Distinct offset residues with nonempty root sets instead have
negative covariance. These probabilities are on a finite residue
space; they are not probabilities of prime events on the integers.

\begin{theorem}[Joint prime-sieve density]\label{thm:joint}
For a finite offset family $\mathcal A$, let
\[
 \beta_{\mathcal A}(p)=\left|\bigcup_{a\in\mathcal A}R_a(p)\right|
 =\sum_{c\in\mathcal A\bmod p}\rho_c(p),
\]
where the sum uses one representative of each distinct offset residue.
The number of classes modulo $P$ surviving for every offset is
\[
 N_{\mathcal A}(P)=\prod_{p\mid P}(p-\beta_{\mathcal A}(p)),
 \qquad
 \delta_{\mathcal A}(P)=\frac{N_{\mathcal A}(P)}P
 =\prod_{p\mid P}\left(1-\frac{\beta_{\mathcal A}(p)}p\right).
\]
\end{theorem}
\begin{proof}
Lemma~\ref{lem:overlap} makes the root sets of distinct offset residues
disjoint, proving the sum formula. There are $p-\beta_{\mathcal A}(p)$
good local classes. CRT identifies the global good classes with their
Cartesian product. Divide the resulting count by $P$.
\end{proof}

For two offsets, the bad count is $\rho_a(p)$ when $p\mid a-b$,
and $\rho_a(p)+\rho_b(p)$ otherwise. Thus for $\{5,11\}$ at $3$,
the joint good proportion is $1/3$, whereas multiplying the separate
good proportions would incorrectly give $1/9$.
Across distinct primes the product is exact because CRT supplies the
bijection. Within one prime, the root union must be counted first.

\section{Local squarefreeness modulo prime squares}
The prime sieve removes every multiple of a selected prime.
To test squarefreeness locally, we should instead remove only the
multiples of its square. This changes the local problem: a root
modulo $p$ need not lift, and a repeated root can have many lifts.

\begin{proposition}[Prime-square root count]\label{prop:square-roots}
For an odd prime $p$ and an integer $a$,
\[
 \rho_a(p^2)=
 \begin{cases}
 1+\leg{-a}{p},&p\nmid a,\\
 0,&p\mid a,\quad p^2\nmid a,\\
 p,&p^2\mid a.
 \end{cases}
\]
At $p=2$, the count $\rho_a(4)$ is $4$ if $4\mid a$ and $0$ otherwise.
\end{proposition}
\begin{proof}
At an odd prime, $x=2m$ is bijective modulo $p^2$, so the equation
is $x^2\equiv-a\pmod{p^2}$.

Suppose $p\nmid a$. Each root $x_0$ modulo $p$ is nonzero.
Choose its representative $0\le x_0<p$ and write
$x_0^2+a=pc$. Every lift has the form $x=x_0+pt$ with $0\le t<p$.
Expanding and dividing by $p$ shows that it is a root modulo $p^2$
exactly when
\[
 c+2x_0t\equiv0\pmod p.
\]
The coefficient $2x_0$ is invertible, so there is exactly one $t$.
Every root modulo $p^2$ reduces to a root modulo $p$, and different
roots have different reductions. The count is therefore unchanged.

If $p\mid a$, reduction modulo $p$ forces $p\mid x$.
Thus $p^2\mid x^2$, and the original congruence can hold only if
$p^2\mid a$. This proves the zero count when $p$ divides $a$ just once.
When $p^2\mid a$, the roots are precisely $x=pk$ modulo $p^2$,
with $0\le k<p$, giving $p$ roots.

Finally, $4m^2+a\equiv a\pmod4$ for all four classes of $m$.
\end{proof}

For example, $G_3$ has one root modulo $3$ but none modulo $9$;
$G_9$ has the three roots $0,3,6$ modulo $9$.
The zero offset falls in the last odd-prime case, and at $2$ every
class is bad. An even offset such as $a=2$ has every class bad
modulo $2$ but no bad class modulo $4$.

\begin{theorem}[Joint finite prime-square sieve]\label{thm:square-sieve}
For a finite family of integer offsets, define
\[
 \beta^{(2)}_{\mathcal A}(p)
 =\left|\bigcup_{a\in\mathcal A}R_a(p^2)\right|
 =\sum_{c\in\mathcal A\bmod p^2}\rho_c(p^2).
\]
The number of classes modulo $M$ for which every $G_a$ avoids
every selected prime square is
\[
 N^{(2)}_{\mathcal A}(M)
 =\prod_{p\in\Pset}(p^2-\beta^{(2)}_{\mathcal A}(p)),
 \qquad
 \delta^{(2)}_{\mathcal A}(M)
 =\frac{N^{(2)}_{\mathcal A}(M)}M
 =\prod_{p\in\Pset}\left(1-\frac{\beta^{(2)}_{\mathcal A}(p)}{p^2}\right).
\]
In particular, the density is positive exactly when
$\beta^{(2)}_{\mathcal A}(p)<p^2$ for every selected prime.
\end{theorem}
\begin{proof}
Apply Lemma~\ref{lem:overlap} with $d=p^2$. Offsets congruent modulo
$p^2$ give identical bad sets, while distinct residues give disjoint
bad sets. This proves the local sum, including the case $p=2$.
The moduli $p^2$ for different primes are pairwise coprime,
so CRT multiplies the numbers of good classes. Each factor is
nonnegative, and a finite product is positive exactly when each
factor is positive.
\end{proof}

For one offset the density is $\prod_{p\in\Pset}(1-\rho_a(p^2)/p^2)$.
Both empty-set conventions agree with the theorem: with no primes
the period is $1$ and its one class survives; with no offsets every
class survives. Repeating an offset adds no restriction. Offsets
congruent modulo $p$ can nevertheless be distinct modulo $p^2$:
$5$ and $11$ coincide modulo $3$, but their bad sets modulo $9$
are disjoint, each with two roots.

The test here is $p^2\nmid G_a(m)$, not
$\gcd(G_a(m),M)=1$. The latter would forbid $p$ itself and would
simply repeat the prime sieve. For instance, $G_3(3)=39$ survives
the square-of-$3$ test despite being divisible by $3$.

\section{The opposite-sign pair}
The family $\{q,-q\}$ permits an explicit comparison of the two sieves.
The prime $q$ behaves differently from every other prime.

\begin{theorem}[No local obstruction for either sieve]\label{cor:signs}
Let $q$ be an odd prime and $\mathcal A=\{q,-q\}$. Its local bad counts are
\[
\begin{array}{c|cc}
 p & \beta_{\mathcal A}(p) & \beta^{(2)}_{\mathcal A}(p)\\ \hline
 p=2 &0&0\\
 p=q &1&0\\
 p\ne q,\ p\equiv3\pmod4 &2&2\\
 p\ne q,\ p\equiv1\pmod4 &2+2\leg{q}{p}&2+2\leg{q}{p}
\end{array}
\]
Thus $\beta_{\mathcal A}(p)<p$ and
$\beta^{(2)}_{\mathcal A}(p)<p^2$ for every prime. Both finite
survivor densities are positive for every finite prime set.
\end{theorem}
\begin{proof}
At $2$, both offsets are odd, so there are no roots modulo $2$ or $4$.
At $q$, both offsets give the root $0$ modulo $q$; neither offset
is divisible by $q^2$, so neither polynomial has a root modulo $q^2$.

At an odd prime $p\ne q$, the offsets are distinct modulo $p$ and
modulo $p^2$, since $p\nmid2q$. Their root sets are disjoint at
both moduli. Also $p\nmid q$, so the individual counts are unchanged
under lifting. At either modulus the union count is
\[
 2+\leg{q}{p}+\leg{-q}{p}.
\]
The identities $\leg{-q}{p}=\leg{-1}{p}\leg{q}{p}$ and
$\leg{-1}{p}=(-1)^{(p-1)/2}$ give the table.

These identities follow from Euler's criterion, which can be seen
directly here. The squaring map on $\mathbb F_p^\times$ has kernel
$\{1,-1\}$ and hence $(p-1)/2$ distinct square values. By Fermat's
identity these are all roots of $x^{(p-1)/2}-1$; every other nonzero
element is a root of $x^{(p-1)/2}+1$. Thus
$t^{(p-1)/2}$ represents the Legendre symbol, and exponentiation
respects multiplication. Substituting $t=-1$ gives its sign.

For $p\equiv3\pmod4$, $2<p$, including $p=3$.
For $p\equiv1\pmod4$, $p\ge5$ and the count is at most $4<p$.
The cases $2,q$ are also strict. The square counts are at most $4$
at the remaining odd primes, less than $p^2$.
Theorems~\ref{thm:joint} and \ref{thm:square-sieve} now give positivity.
\end{proof}

In particular, the prime-square density can be written
\[
 \delta^{(2)}_{\{q,-q\}}(M)
 =\prod_{\substack{p\in\Pset\\p\ne2,q}}
 \left(1-\frac{2+\leg{q}{p}+\leg{-q}{p}}{p^2}\right).
\]
No factor is needed at $2$ or $q$. At $q$, prime-square survival
allows all centers, whereas prime survival removes $m\equiv0\pmod q$.

\begin{corollary}[A wider prime-square admissible family]\label{cor:family}
Every finite family of nonzero squarefree integer offsets has positive
density under every finite prime-square sieve. In particular, this
holds for finite families of signed prime offsets.
\end{corollary}
\begin{proof}
No prime square divides any offset $a$ in such a family. Consequently,
$m=0$ is good modulo $p^2$ for every offset and every prime $p$.
Each local good set is nonempty, so Theorem~\ref{thm:square-sieve} applies.
\end{proof}

Indeed, every multiple of $P$ is such a prime-square survivor,
since $G_a(m)\equiv a\pmod{p^2}$ for each selected $p$.
This does not make larger signed-prime families
admissible for the prime sieve: $\{3,5,7\}$ already obstructs it at $3$.

Each positive finite density supplies a surviving residue class,
which repeats indefinitely. This guarantees centers surviving any
one fixed finite sieve. It says nothing about surviving all prime
squares simultaneously or producing simultaneous primes.

\section{Exact implementation and interval evaluation}
The code constructs local roots, takes the union over offsets,
and combines the good classes by CRT. For prime squares, a
nondegenerate root $r$ of $G_a$ modulo $p$ is lifted directly in
the $m$ variable:
\[
 c=\frac{4r^2+a}{p},\qquad
 t\equiv-c(8r)^{-1}\pmod p,\qquad m=r+pt.
\]
Here $8r$ is invertible when $p$ is odd and $p\nmid a$.
This follows by expanding $G_a(r+pt)$; it is the same one-step
argument as Proposition~\ref{prop:square-roots}. The singular cases
and modulus $4$ are constructed separately.

Suppose $r$ is good modulo the current period $D$ and $s$ is good
modulo a coprime next modulus $h$, either $p$ or $p^2$. The combined
canonical class is
\[
 r+D\big((s-r)D^{-1}\bmod h\big).
\]
The parenthesized residue lies in $\{0,\ldots,h-1\}$, so the result
is in $[0,Dh)$ and has the required two residues. Each pair gives
one distinct class, implementing the CRT bijection.

\begin{proposition}[Exact interval evaluation]\label{prop:exact}
Let $D\ge1$ be a period and $\mathcal R\subseteq\{0,\ldots,D-1\}$
its survivor classes, with $N=|\mathcal R|$. Then
\[
 S(L,U)=\sum_{r\in\mathcal R}
 \left(\left\lfloor\frac{U-r}D\right\rfloor-
       \left\lfloor\frac{L-r}D\right\rfloor\right).
\]
If $H=tD+h$, $0\le h<D$, then
\[
 S(L,U)=tN+T,\qquad 0\le T\le\min(h,N).
\]
In particular $|S(L,U)-HN/D|\le N$, with equality of the count
and main term whenever $D\mid H$.
\end{proposition}
\begin{proof}
Sum Lemma~\ref{lem:interval} over $\mathcal R$. Partition the interval
into $t$ consecutive blocks of $D$ integers and a remainder of $h$
integers. Every full block contains each residue once; the remainder
contains no repeated residue and at most $\min(h,N)$ good ones.
Both $T$ and $hN/D$ lie in $[0,N]$, giving the error bound.
If $h=0$, both are zero.
\end{proof}

Take $D=P$ for the prime sieve and $D=M$ for the prime-square sieve.
For one shift, either this bound or the inclusion-exclusion bound
can be smaller. The program uses the exact floor sum rather than
rounding a density times the interval length.

The reusable routines are in \texttt{src/arithmetic.py},
\texttt{src/sieve.py}, and \texttt{src/squarefree.py}.
Moduli are checked to be distinct primes. Empty prime sets use
the period $1$. Local counts and densities use integers and rational
arithmetic. Direct divisibility enumeration checks the constructed
classes and interval counts independently. These checks verify
the implementation; the proofs establish the formulas.
Explicit class storage grows with the period and is intended for
modest products. The repository's historical name
\texttt{quadratic-sieve-lab} does not refer to an implementation of
the classical quadratic-sieve factorization algorithm.

\section{Divisor compatibility and witness moments}
The local sieve fixes offsets. For the separate witness problem,
two elementary tools help describe collisions and occupancy.

\begin{proposition}[Divisor compatibility]\label{prop:quotient}
Let $d,e$ be positive integers with $d\mid G_a(m)$. Write
$g=\gcd(d,e)$, $d=gu$, $e=gv$, and $k=G_a(m)/d$. Then
\[
 e\mid G_b(m)\iff
 g\mid a-b\ \hbox{ and }\ uk\equiv(a-b)/g\pmod v.
\]
In particular, simultaneous divisibility requires $\gcd(d,e)\mid a-b$.
\end{proposition}
\begin{proof}
Subtraction gives $G_b(m)=guk-(a-b)$. If $e$ divides this value,
then $g\mid a-b$. Under that condition, dividing by $g$ leaves
exactly $v\mid uk-(a-b)/g$. Reverse the steps for sufficiency.
\end{proof}

The gcd condition alone is insufficient: with $a=3,b=5,m=1,d=7,e=5$,
we have $d\mid G_a(1)$ and $\gcd(d,e)=1\mid a-b$, but
$5\nmid G_b(1)=9$. Equal signs involve $q-r$ up to sign;
opposite signs involve $q+r$ up to sign. No squarefreeness
assumption on $d,e$ is needed.

Let $\mathcal C$ be a finite set of $H$ centers. For one sign and
witness class, let $\nu(n)$ count successful offsets and define
\[
 M_1=\sum_{n\in\mathcal C}\nu(n),\quad
 M_2=\sum_{n\in\mathcal C}\nu(n)^2,\quad
 F_2=\sum_{n\in\mathcal C}\nu(n)(\nu(n)-1),\quad
 S=\#\{n\in\mathcal C:\nu(n)>0\}.
\]
The factorial moment $F_2$ counts ordered pairs of distinct successful
offsets at the same center.

\begin{proposition}[Moments and occupancy]\label{prop:moments}
Always $M_2=M_1+F_2$ and $S\le\min(H,M_1)$. If $M_2>0$, then
\[
 \frac{M_1^2}{M_2}\le S,\qquad
 0\le M_1-S\le\frac{F_2}{2}.
\]
If $M_2=0$, all counts vanish; the lower bound is then interpreted as $0$.
\end{proposition}
\begin{proof}
For a nonnegative integer $j$, $j^2=j+j(j-1)$.
Every occupied center contributes at least one to $M_1$.
Cauchy--Schwarz on the occupied centers gives $M_1^2\le SM_2$.
Finally, $0\le j-\ind_{j>0}\le j(j-1)/2$: both sides vanish
at $j=0$, and for $j\ge1$ this is $j-1\le j(j-1)/2$.
Sum these inequalities. If $M_2=0$, every $\nu(n)^2$ is zero.
\end{proof}

If one separately proved $M_1\ge cHK$ and $M_2\le CH(K+K^2)$
with $c,C,K>0$, this would imply
\[
 S\ge\frac{c^2}{C}H\frac{K}{1+K}.
\]
The moment hypotheses at unbounded scales are not proved here.

For the sets $A_+,A_-\subseteq\mathcal C$ occupied on each sign,
inclusion-exclusion gives
\[
 |A_+\cap A_-|\ge\max(0,|A_+|+|A_-|-H).
\]
Two disjoint sets can each occupy $40\%$ of the centers.
Positive separate occupancy need not force overlap. Even when
$A_+\cap A_-$ is nonempty, the two witnesses can use different offsets.

\section{Exact finite experiments}
\subsection{Fixed-offset prime and prime-square sieves}
For $\Pset=\{2,3,5,7\}$, the prime period is $P=210$ and the
prime-square period is $M=44{,}100$. Both experiments construct
CRT classes, check a complete period by direct enumeration,
and count the interval $500<m\le1000$.
The following tables are exact finite counts.

\begin{center}
\begin{tabular}{lrrr}
\toprule
Offsets & Classes mod $210$ & Density & Interval count\\
\midrule
$\{3\}$ &100&$10/21$&237\\
$\{-3\}$ &140&$2/3$&333\\
$\{3,-3\}$ &100&$10/21$&237\\
$\{5\}$ &40&$4/21$&95\\
$\{-5\}$ &168&$4/5$&400\\
$\{5,-5\}$ &40&$4/21$&95\\
$\{5,11\}$ &20&$2/21$&47\\
\bottomrule
\end{tabular}
\end{center}

\begin{center}
\begin{tabular}{lrrr}
\toprule
Offsets & Classes mod $44{,}100$ & Density & Interval count\\
\midrule
$\{3\}$ &42,300&$47/49$&480\\
$\{-3\}$ &44,100&$1$&500\\
$\{3,-3\}$ &42,300&$47/49$&480\\
$\{5\}$ &32,900&$47/63$&372\\
$\{-5\}$ &44,100&$1$&500\\
$\{5,-5\}$ &32,900&$47/63$&372\\
$\{5,11\}$ &21,620&$1081/2205$&246\\
$\{9\}$ &27,048&$46/75$&305\\
\bottomrule
\end{tabular}
\end{center}

For $\{5,-5\}$, the prime-square bad counts at $2,3,5,7$ are
$0,2,0,2$, giving $(1-2/9)(1-2/49)=47/63$.
For $\{5,11\}$, the bad count at $3^2$ is $4$, despite the coincident
root sets modulo $3$. These examples show both the change from prime
to prime-square conditions and the need to group offsets at the
correct modulus. Density $1$ for $\{-5\}$ refers only to these
four selected squares.

The prime-square experiment also includes offsets $0,2,4$, repeated
offsets and an empty family. Tests cover empty prime sets, singular
roots, negative intervals and invalid moduli. The interval length is
$500$, not a multiple of either period, so an interval count need not
equal $500$ times its period density.

\subsection{Witnesses with varying prime offsets}
The original experiment tests the $500$ even centers $1000<n\le2000$
and the $13$ odd primes $3\le q\le43$ on both signs, giving
$500\cdot13\cdot2=13{,}000$ values. Trial division fully factors each
positive value. Every factor is checked for primality and the product
is checked against the input. Failures are retained alongside successes.

\begin{center}
\begin{tabular}{rlrrr}
\toprule
$R$ & Witness class & Upper occupied & Lower occupied & Both sides\\
\midrule
7 & Prime &216&236&104\\
7 & Rough squarefree $P_2$ &454&461&421\\
19 & Prime &329&380&249\\
19 & Rough squarefree $P_2$ &497&498&495\\
43 & Prime &439&448&390\\
43 & Rough squarefree $P_2$ &499&500&499\\
\bottomrule
\end{tabular}
\end{center}

The last column allows different offsets on the two signs.
At $R=7$, upper prime witnesses have $M_1=246$, $M_2=306$,
$F_2=60$, and $S=216$. Thus the lower bound is
$246^2/306=3362/17$, and $M_1-S=30=F_2/2$.
The saved summary includes moments for all twelve sign/class/cutoff
combinations. Rough squarefree $P_2$ occupancy is much higher in this
window; it cannot be relabeled as prime occupancy.

\subsection{Reproduction}
From the repository root, with Python 3.10 or later and no third-party
packages, run
\begin{verbatim}
python -m unittest discover -s tests -v
python -m experiments.witness_counts
python -m experiments.local_sieve
python -m experiments.squarefree_sieve
\end{verbatim}
The witness experiment writes \texttt{results/witnesses.csv} and
\texttt{results/summary.json}. The two sieve experiments each write
a JSON record and a Markdown table in \texttt{results/}.
The checked-in default results correspond to the tables above.
Each module accepts \texttt{--help} and \texttt{--output}; the latter
keeps alternative runs separate. The legacy entry point
\texttt{experiment.py} still runs the witness experiment.

Tests compare every offset residue modulo $p^2$ for primes through
$19$ with brute-force roots, verify overlaps, and compare small CRT
periods and intervals with direct divisibility tests. The original
$272$ root/count checks and $47{,}670$ quotient checks are retained.
These ranges describe software verification, not the scope of the proofs.

\section{Limitations and status}
A finite sieve survivor need not be prime or globally squarefree.
For example, $G_5(2)=21$ avoids the primes $2,5,11$ but is composite.
Also $G_9(2)=25$ avoids the prime squares $4,9$ but is not squarefree.
Conversely, a prime value equal to a sieving prime is removed:
$G_{-11}(2)=5$ fails the prime sieve containing $5$.

For a positive value $v>1$, checking every prime up to $\sqrt v$
does certify primality if no divisor is found: a composite value has
a prime divisor at most $\sqrt v$. Checking prime-square divisibility
for every prime up to $\sqrt v$ likewise certifies squarefreeness.
Both are bounded tests for a specified value, rather than consequences
of one fixed finite prime set as $m$ grows.

The CRT densities are exact for periodic local survivor sets.
They are not densities of prime values or of values avoiding every
prime square. Positivity for every finite prime set does not by itself
justify passing to an infinite set of restrictions. The opposite-sign
theorem proves local admissibility and supplies no global
simultaneous-squarefreeness or simultaneous-prime theorem.

The witness experiment tests exact factorization, squarefreeness,
the $P_2$ condition and roughness. A finite prime-square sieve alone
does not supply those properties. Its tables concern the stated window
and cutoffs, and the moment inequalities require separately proved
moment estimates for any asymptotic deduction. No estimate uniform
in a growing sieve period is claimed.

The project uses elementary tools: quadratic residues, CRT, finite
inclusion-exclusion, and Cauchy--Schwarz. The development here is their
combination for several shifted squares, the extension to prime-square
local conditions, and an exact implementation. Proofs of the stated
results are included.

\section{Further questions}
\paragraph{1. Higher prime powers.}
Determine $\rho_a(p^k)$ for $k\ge3$, separating the possible orders
of divisibility of $a$. The nondegenerate one-step lift can be iterated;
the repeated-root cases require a separate count. Compare finite
power-free sieves and the exact periods of their survivor sets.

\paragraph{2. Larger families and prime obstructions.}
Classify minimal signed-prime-offset families with
$\beta_{\mathcal A}(p)=p$, beginning at $3$ and $5$.
Such families obstruct the prime sieve even though
Corollary~\ref{cor:family} excludes prime-square obstructions.
The criterion is whether the offset residues cover all values
$-4m^2$ modulo $p$; the challenge is to give a useful classification.

\paragraph{3. Discrepancy under translation.}
For one fixed survivor set of period $D$ and a fixed length $H$,
determine the maximum and minimum of $S(L,L+H)-HN/D$ as
$L$ runs modulo $D$. The remainder count in
Proposition~\ref{prop:exact} gives a starting bound.
Can the arrangement of good classes improve it for these sieves?
This is a finite combinatorial question for either kind of sieve.
\end{document}
